\section{Synthesis: answers to the research questions}
\label{sec:synthesis}

\paragraph{RQ1---semantic measurement.}
The controlled retrieval order responds to width and effective data, and its
risk decomposition is auditable. This supports the measurement on the
registered synthetic task. It does not prove that the same order survives
natural language.

\paragraph{RQ2---finite-size boundary.}
The answer is negative on the public reference grid. There is no interior
susceptibility peak, no Binder crossing, no attention specialization, and no
prospective largest-size prediction. The correct outcome is finite-width
learning response, not phase transition.

\paragraph{RQ3---mechanism.}
Matched ablations support a bounded feed-forward contribution. The frozen suite
also separates the registered MLP variants in task order. The result is
task-conditional and does not establish an architectural phase split.

\paragraph{RQ4---optimizer and scale path.}
Frozen slices show substantial optimizer-conditioned risk differences, and
uncapped data paths outperform the linear data path at the registered endpoint.
Because compute, update geometry, and all hyperparameters are not fully matched,
the result remains a contrast within the atlas rather than a global method
ranking.

\paragraph{RQ5---predictive transport.}
The controlled continuation simulator verifies that calibration, censoring,
baselines, and disjoint holdouts can be executed. Its generated response law
makes near-exact prediction unsurprising; it is therefore not used as empirical
evidence for transport in trained systems. A valid positive answer requires the
dense Transformer experiment to predict an untouched size or realism edge.

\paragraph{RQ6---external scope.}
Natural corpora support tokenizer-independent bits-per-byte measurement on
disjoint ranges. Diffusion, RL, and multi-agent artifacts demonstrate protocol
portability. Learned diffusion endpoints, open RLVR, open-weight large models,
and truth-labeled agent populations remain open. External validity is therefore
partial and explicitly bounded.

\subsection{What the taxonomy changes in practice}

Three distinctions drive these answers. First, nominal control is separated
from effective control, revealing the duplicated data point. Second, causal
intervention evidence is separated from phase evidence, so a real MLP effect
does not manufacture a Binder crossing. Third, artifact coverage is separated
from evidential depth, so a cross-domain simulator does not become a natural
endpoint merely because it shares plotting code.

The result is a coherent claim ledger rather than twenty independent figures.
Each positive statement has a neighboring blocked statement and a next
falsifier. This makes negative evidence cumulative: a future study can change
the status by filling a named missing coordinate rather than restarting the
terminology.

\FloatBarrier
